% ═══════════════════════════════════════════════════════════════════
% ЧАСТЬ V (продолжение). ТЕОРЕМЫ 1–3 СТЕНДА N=15/30
% ═══════════════════════════════════════════════════════════════════

\section{Теорема~1. Соотношения Римана как тождество}\label{sec:riemann}

\begin{definition}[нормированная период-матрица]\label{def:periodmatrix}
Пусть $\mathfrak{B}=\{a_1,\dots,a_g,b_1,\dots,b_g\}$ --- симплектический
базис леммы~\ref{lem:sympl}. Нормируем базис форм
$\omega_1,\dots,\omega_g$ условием $\int_{a_i}\omega_j=\delta_{ij}$ и
положим
\[
  \Omega \;=\; \bigl(\Omega_{ij}\bigr)_{i,j=1}^{g},
  \qquad \Omega_{ij}=\int_{b_i}\omega_j .
\]
\end{definition}

\begin{theorem}[соотношения Римана]\label{th:riemann}
Для период-матрицы определения~\ref{def:periodmatrix} выполнены
\begin{enumerate}[leftmargin=2em, itemsep=0.15em]
\item $\Omega^{\mathsf{T}}=\Omega$ (симметрия);
\item $\im\Omega$ положительно определена: для любого
  $0\ne\lambda\in\RR^g$ имеет место
  $\lambda^{\mathsf{T}}(\im\Omega)\lambda>0$.
\end{enumerate}
\end{theorem}

\begin{proof}
Инструмент --- билинейное тождество Римана. Пусть $F$ ---
фундаментальный $4g$-угольник леммы~\ref{lem:sympl} с рёбрами
$a_i,b_i$, и $\alpha,\beta$ --- замкнутые голоморфные 1-формы. Формула
Стокса, применённая к $\alpha\wedge\beta$ на $F$, с учётом сокращения
противоположных рёбер даёт
\begin{equation}\label{eq:bilinear}
  \iint_F \alpha\wedge\beta
  \;=\;\sum_{i=1}^{g}\Bigl(
    \int_{a_i}\alpha\int_{b_i}\beta
    \;-\;\int_{b_i}\alpha\int_{a_i}\beta\Bigr).
\end{equation}

\emph{Симметрия.} Возьмём $\alpha=\omega_j$, $\beta=\omega_k$ --- оба
голоморфны. На комплексной кривой всякий голоморфный $(2,0)$-объект
обращается в нуль (базисное поле форм типа $(2,0)$ пусто:
$\Omega^2_{C_N}=0$), поэтому $\omega_j\wedge\omega_k=0$, и
\eqref{eq:bilinear} даёт
$\sum_i(\delta_{ij}\Omega_{ik}-\Omega_{ij}\delta_{ik})
=\Omega_{jk}-\Omega_{kj}=0$.

\emph{Положительность.} Возьмём произвольную ненулевую голоморфную
форму $\omega=\sum_k\lambda_k\omega_k$, $\lambda\in\CC^g\setminus
\{0\}$, и применим \eqref{eq:bilinear} к $\alpha=\omega$,
$\beta=\overline{\omega}$. Локально $\omega=g\,dz$, и
$\omega\wedge\overline{\omega}=|g|^2\,dz\wedge d\bar z
=-2i\,|g|^2\,dx\wedge dy$ --- отрицательная вещественная 2-форма, так
что левая часть \eqref{eq:bilinear} равна
$-2i\int_{C_N}|g|^2\,dx\,dy$ --- чисто мнимое число. Правая часть с
учётом $\int_{a_i}\omega=\lambda_i$ и
$\int_{b_i}\omega=(\Omega\lambda^{\mathsf{T}})_i$ равна
\[
  \sum_i\Bigl(\lambda_i\,\overline{(\Omega\lambda^{\mathsf{T}})_i}
  -(\Omega\lambda^{\mathsf{T}})_i\,\overline{\lambda_i}\Bigr)
  \;=\;2i\,\im\!\Bigl(\lambda\,\overline{\Omega}\,\overline{\lambda}^{\mathsf{T}}\Bigr)
  \;=\;-2i\,\lambda^{\mathsf{T}}(\im\Omega)\overline{\lambda}^{\mathsf{T}},
\]
где во втором переходе использованы симметрия $\Omega$ (пункт~1) и
тождество
$\im(\lambda\overline{\Omega}\overline{\lambda}^{\mathsf{T}})
=-\lambda^{\mathsf{T}}(\im\Omega)\overline{\lambda}^{\mathsf{T}}$.
Приравнивая обе части и сокращая $-2i$:
\[
  \lambda^{\mathsf{T}}(\im\Omega)\overline{\lambda}^{\mathsf{T}}
  \;=\;\int_{C_N}|g|^2\,dx\,dy\;>\;0,
\]
и для вещественного вектора $\lambda\ne0$ это даёт утверждение~2
(положительная определённость $\im\Omega$).
\end{proof}

\subsection{Почему здесь не нужен самосопряжённый оператор}\label{sec:nooperator}

Слово «симплектический» в настоящем тексте относится исключительно к
косой форме пересечений $J$ на целочисленной решётке
$H_1(C_N,\ZZ)$ --- объекту чисто топологическому
(лемма~\ref{lem:sympl}); это не фазовое пространство механики, и ни
гамильтониан, ни наблюдаемые, ни спектральная задача в доказательстве
не фигурируют. В теореме~\ref{th:riemann} участвуют ровно два объекта,
и оба построены внутри нашей конструкции: решётка циклов с формой
пересечений (раздел~\ref{sec:cycles}) и интегралы голоморфных форм по
циклам (раздел~\ref{sec:periods}). Соотношения Римана --- не внешняя
гипотеза и не спектральное утверждение, а \emph{тождество, которое
наша собственная период-матрица обязана выполнять}. Именно поэтому
численная проверка на матрицах $91\times182$ и $406\times812$
демонтирована как этап: проверять числом нечего --- доказано
тождественно. Числа остаются только в качестве одноразовой
санити-строки вычислительного протокола, вне текста. Единственная
содержательная причина сохранить сам блок --- форма $J$ является
входом модуля индексов решётки Ходжа (раздел~\ref{sec:indices}):
без поляризации индексы не определяются. Этот вход получен из первых
принципов и стоит двух страниц.

\section{Теорема~2. $\Ga$-нормировка и дробные доли $\pi$}\label{sec:normalization}

\begin{definition}[$\Ga$-нормировка стенда]\label{def:grossnorm}
Для характера $(a,b)\in T_N$ и периода
$P\in\frac{\Omega_{a,b}}{N}\ZZ[\zeta_N]$ (лемма~\ref{lem:closedform})
\emph{нормированным периодом} называется
\[
  \widehat{P} \;=\; \frac{P}{\Omega_{a,b}}
  \;\in\;\frac1N\,\ZZ[\zeta_N].
\]
Трансцендентный множитель полностью поглощается величиной
$\Omega_{a,b}=\Bfun(a/N,b/N)$; нормированные периоды ---
алгебраические числа.
\end{definition}

\begin{theorem}[$\Ga$-нормировка и дробные доли $\pi$]\label{th:normalization}
Пусть $N\in\{15,30\}$. Тогда:
\begin{enumerate}[leftmargin=2em, itemsep=0.15em]
\item \emph{Алгебраичность.} Для всякого периода $P$ собственной
  формы $(a,b)$ нормировка $\widehat P=P/\Omega_{a,b}$ лежит в
  $\frac1N\ZZ[\zeta_N]$: фазовая решётка стенда имеет знаменатель
  $N$, то есть $15$ и $30$ соответственно.
\item \emph{Редукция трансцендентной части.} Каждое $\Omega_{a,b}$
  записывается через значения $\Ga$ в канонических позициях и
  множители вида $\pi/\sin(\pi k/N)$, $1\le k\le N-1$, согласно
  лемме~\ref{lem:relations}; в случае $a+b=N$ имеем чистую дробную
  долю $\Omega_{a,b}=\pi/\sin(\pi a/N)$.
\item \emph{Отсутствие $4\pi^2$ в нормировке.} В замкнутой
  форме~\eqref{eq:closedform} число $\pi$ входит единственным
  образом --- через множители $\pi/\sin(\pi k/N)$ с $k/N=k/15$ или
  $k/30$ и через формулу умножения~\eqref{eq:gauss}; никакая степень
  вида $4\pi^2$ не возникает ни в нормировке, ни в промежуточных
  соотношениях. (Универсальный Tate-фактор $4\pi^2$ живёт в спектре
  решёток --- сертификат~G, теорема~\ref{th:certG} --- и в нормировку
  периодов не входит.)
\item \emph{Замкнутые радикальные константы стенда.} Обобщённая
  константа $b_{Ch}(n)=1-\cos(2\pi/n)$ на уровнях стенда принимает
  точные радикальные значения:
  \begin{align}
    b_{Ch}(15) &= 1-\cos\frac{2\pi}{15}
      \;=\;1-\frac{1+\sqrt5+\sqrt3\,\sqrt{10-2\sqrt5}}{8},
      \label{eq:b15}\\
    b_{Ch}(30) &= 1-\cos\frac{\pi}{15}
      \;=\;1-\frac{\sqrt{30+6\sqrt5}+\sqrt5-1}{8}.
      \label{eq:b30}
  \end{align}
\end{enumerate}
\end{theorem}

\begin{proof}
Пункт~1 --- немедленное следствие леммы~\ref{lem:closedform}: образ
период-функционала лежит в $\frac{\Omega}{N}\ZZ[\zeta_N]$, деление
на $\Omega$ даёт $\frac1N\ZZ[\zeta_N]$. Пункты~2 и~3 --- применение
леммы~\ref{lem:relations}: все используемые преобразования значений
$\Ga(k/N)$ исчерпываются отражением, рекуррентностью и умножением;
отражение вводит $\pi$ исключительно в виде $\pi/\sin(\pi k/N)$, а
умножение --- в виде $(2\pi)^{(m-1)/2}$ при целых $m$; ни та, ни
другая операция не порождает нормировочного множителя $4\pi^2$,
поскольку нормировка стенда строится из самих периодов. Пункт~4:
точные значения --- косинусы половинных аргументов:
$\cos24^\circ=\cos(60^\circ-36^\circ)$ с
$\cos36^\circ=\frac{1+\sqrt5}{4}$,
$\sin36^\circ=\frac{\sqrt{10-2\sqrt5}}{4}$ даёт~\eqref{eq:b15};
аналогично $\cos12^\circ=\cos(30^\circ-18^\circ)$ с
$\cos18^\circ=\frac{\sqrt{10+2\sqrt5}}{4}$,
$\sin18^\circ=\frac{\sqrt5-1}{4}$ даёт~\eqref{eq:b30}
(сверка численно: $b_{Ch}(15)\approx0.086\,454$,
$b_{Ch}(30)\approx0.021\,854$).
\end{proof}

\begin{remark}[о смысле «дробных долей $\pi$»]
Доля $\pi$ в нормировке всегда рациональна со знаменателем $N$ ---
множитель $\pi/\sin(\pi k/N)$ несёт угол $k\pi/N$ (при $N=15,30$ это
доли $\pi/15$ и $\pi/30$), а фазовая решётка $\frac1N\ZZ[\zeta_N]$
несёт знаменатель $N$ в алгебре. Это структурное утверждение: оно
верно для всех блоков сразу и не зависит от выбора выборки
характеров. Лестница уровней $\pi/7$, $\pi/15$, $\pi/30$
(сертификаты~G и~H) показывает, что все ступени программы --- ступени
одной лестницы $\mu_N$-квантов.
\end{remark}

\section{Теорема~3. Поляризация и индексы решётки Ходжа}\label{sec:indices}

\begin{definition}[билинейная форма на когомологиях и поляризация]
\label{def:polarization}
Для замкнутых 1-форм $\alpha,\beta$ на $C_N$ положим
$\widetilde E([\alpha],[\beta])=\iint_{C_N}\alpha\wedge\beta$; по
Стоксу форма $\widetilde E$ корректно определена на
$H^1_{\mathrm{dR}}(C_N)$ и кососимметрична. Поляризационной формой на
$H_1(C_N,\ZZ)$ называется форма пересечений $E:=J$; по двойственности
Пуанкаре $\widetilde E(\delta_u,\delta_v)=E(u,v)$, где
$\delta_u,\delta_v$ --- классы, двойственные Пуанкаре к циклам.
\end{definition}

\begin{theorem}[$J$ --- поляризация структуры Ходжа веса $1$]
\label{th:polarization}
Относительно разложения Ходжа
$H^1(C_N,\CC)=H^{1,0}\oplus H^{0,1}$, где
$H^{1,0}=\operatorname{span}\{\omega_1,\dots,\omega_g\}$ ---
пространство голоморфных форм леммы~\ref{lem:forms}, выполнены
соотношения Римана--Ходжа:
\begin{enumerate}[leftmargin=2em, itemsep=0.15em]
\item $\widetilde E(\omega,\omega')=0$ для всех
  $\omega,\omega'\in H^{1,0}$;
\item $-\dfrac{1}{2i}\,\widetilde E(\omega,\overline{\omega})
  =\displaystyle\int_{C_N}|g|^2\,dx\,dy>0$ для всякой ненулевой
  $\omega=g\,dz\in H^{1,0}$.
\end{enumerate}
Следовательно, $E=J$ --- поляризация веса $1$ на решётке
$H_1(C_N,\ZZ)$, и якобиан $\Jac(C_N)=H^1/H_1$ вместе с этой
поляризацией является главным поляризованным абелевым многообразием
(матрица поляризации в симплектическом базисе --- $J$).
\end{theorem}

\begin{proof}
(i) На комплексной кривой нет ненулевых голоморфных 2-форм:
$\Omega^2_{C_N}=0$, поэтому $\omega\wedge\omega'=0$ как голоморфный
$(2,0)$-объект, и интеграл равен нулю. Эквивалентно, это в точности
первое соотношение Римана из теоремы~\ref{th:riemann}, прочитанное на
билинейном тождестве~\eqref{eq:bilinear}. (ii) Вычисление из
доказательства теоремы~\ref{th:riemann}:
$\omega\wedge\overline{\omega}=|g|^2\,dz\wedge d\bar z
=-2i\,|g|^2\,dx\wedge dy$, откуда
$-\frac{1}{2i}\widetilde E(\omega,\overline{\omega})
=\int_{C_N}|g|^2\,dx\,dy>0$ при $\omega\ne0$ (подынтегральная функция
непрерывна, не равна тождественно нулю и неотрицательна).
Соотношения (i)--(ii) суть в точности условия Римана для того, чтобы
$\im\Omega\succ0$ (теорема~\ref{th:riemann}, пункт~2) --- то есть
чтобы период-матрица задавала главное поляризованное абелево
многообразие; унимодулярность $J$ (лемма~\ref{lem:sympl}) даёт
главность поляризации.
\end{proof}

\begin{definition}[индексы решётки Ходжа по блокам]\label{def:indices}
Разложение на собственные характеры (лемма~\ref{lem:forms}) определяет
по каждому проводнику $d\mid N$ решётку
\[
  \Lambda_d \;=\; \Bigl\{\gamma\in H_1(C_N,\ZZ):\ \gamma
  \text{ лежит в изотипной компоненте характеров }(a,b),\
  d(a,b)=d\Bigr\},
\]
её насыщение $\Lambda_d^{\mathrm{sat}}=\Lambda_d\otimes\QQ\cap
H_1(C_N,\ZZ)$ и \emph{индекс блока}
$\Ind(\Lambda_d)=[\Lambda_d^{\mathrm{sat}}:\Lambda_d]$.
Данные блока: ранг, сигнатура ограничения поляризации
$E|_{\Lambda_d}$, определитель $E|_{\Lambda_d}$ и индекс.
\end{definition}

\begin{theorem}[вычислимость индексов из $\Ga$-данных]\label{th:indexcomp}
Для стенда $N\in\{15,30\}$ все величины определения~\ref{def:indices}
конечны и вычисляются алгоритмически из замкнутых
форм~\eqref{eq:closedform}: насыщение находится целочисленным
приведением решётки блока (нормальная форма Эрмита---Смита),
сигнатура и определитель ограничения поляризации --- из комбинаторно
известной матрицы пересечений (лемма~\ref{lem:sympl}) с фильтрацией
по характерам, индекс --- как отношение определителей. Число
собственных форм по блокам даётся переписью $h_d$
(приложение~\ref{app:census}): для $N=15$ --- три блока с
$1+6+84=91$; для $N=30$ --- шесть блоков с $1+6+9+30+84+276=406$.
\end{theorem}

\begin{proof}
Конечность: $H_1(C_N,\ZZ)$ --- свободная абелева группа конечного
ранга $2g$, подрешётки конечного ранга имеют конечный индекс в
насыщении. Вычислимость: изотипная компонента характера $(a,b)$
определена над $\QQ(\zeta_N)$; её пересечение с целочисленной
решёткой --- целочисленная подрешётка, базис которой находится
алгоритмом Эрмита---Смита по целочисленному базису циклов и таблице
действия $\mu_N\times\mu_N$ (комбинаторика листов
леммы~\ref{lem:gridgraph}). Ограничение поляризации вычисляется по
теореме~\ref{th:polarization}: $E=J$, а матрица $J$ в базисе циклов
известна комбинаторно. Перепись $h_d$ ---
приложение~\ref{app:census}; суммы $h_d$ совпадают с родом (проверено
точной арифметикой, V1 протокола).
\end{proof}

\begin{statusbox}[стык с сертификатом H]
Сертификат~H (раздел~\ref{sec:certH}) реализует вычислимость
теоремы~\ref{th:indexcomp} в полной мере: тип поляризации по блокам
получен в нормальной форме Смита $(1,1,5,5,15,15,15,15)$,
дискриминант $3^4\cdot5^6=1265625=1125^2$, полная $\Ga$-нормировка
объёма $\mathrm{vol}_h=1125$. Таким образом, индексы решётки Ходжа
стенда $N=15/30$ не только вычислимы (теорема~\ref{th:indexcomp}),
но и вычислены с точной целочисленной сертификацией.
\end{statusbox}

\subsection{Стык с программой: тройка $(n,B,\lambda_0)$ и лестница}
\label{sec:repocycle}

Геометрия поставляет только тройку $(n,B,\lambda_0)$ --- порядок
элемента, целочисленный топологический индекс и спектральный порог;
поправки имеют вид скелета
\[
  \Delta(n,B;\lambda_0) \;=\; \lambda_0
  \;+\;\frac{(\pi/n)^2}{2}\;-\;\frac{(\pi/n)^5}{B},
  \qquad
  b_{Ch}(n)=1-\cos\frac{2\pi}{n}.
\]
Стенд $N=15/30$ поставляет компоненты $n$ и $B$ в крупном масштабе:
роль $n$ играет уровень $N$ (порядок корней единицы; все фазы стенда
$\pi/N$-кратны в смысле теоремы~\ref{th:normalization}), а
целочисленные индексы $B$ --- это ранги и определители
поляризационных решёток $\Lambda_d$ (теорема~\ref{th:indexcomp});
на предыдущих ступенях эту роль играл единственный индекс $b_2=22$,
здесь индексов несколько --- по блокам проводников, и они
вычисляются, а не постулируются. Спектральная компонента $\lambda_0$
на стенде не задействована: $\Ga$-нормировка строится из самих
периодов, и никакая спектральная величина в стенд не входит --- это
осознанное разделение ролей между ступенями лестницы, при котором
каждая ступень развивает свои сильные стороны.
